\chapter{Zarourati Backstepping Controller }
\makeatletter\@mkboth{}{Appendix}\makeatother
\label{appen:derivations_tian_oscillator}

The choice of the oscillatory term given by equation \eqref{eq:oscillator} is not immediately obvious.  The oscillatory property of this term can by understood by the time derivative of the term $\vec{e}_d^T\vec{e}_d = \norm{\vec{e}_d}^2$
\begin{align*}
    \frac{d \vec{e}_d^T \vec{e}_d}{dt} &= \vec{e}_d^T \Big[\frac{\dot\xi}{\xi}\vec{e}_d + \Lambda\vec{G}_1\vec{e}_d \Big] \\
    &=2\vec{e}_d^T\frac{\dot\xi}{\xi}\vec{e}_d + 2\vec{e}_d^T\Lambda\vec{G}_1\vec{e}_d \\
    &=2\vec{e}_d^T\frac{\dot\xi}{\xi}\vec{e}_d \numberthis \label{eq:oscillator_norm_ode}
\end{align*}
given $\vec{G}_1 = \begin{bmatrix}0 & 1 \\ -1 & 0\end{bmatrix}$ is skew-symmetric then it has the property where any vector $\vec{x}$ satisfies $\vec{x}^T\vec{G}_1\vec{x} = 0$. Since and $\Lambda$ is a scalar then the second term of the derivative becomes zero.

Given equation \eqref{eq:oscillator_norm_ode}, let $r(t)=\norm{\vec{e}_d(t)}$, so 

\begin{align*}
    \frac{d}{dt}r^2 &= 2\frac{\dot{\xi}}{\xi}r^2 \\
    \therefore \ \frac{dr}{dt}2r &= 2\frac{\dot{\xi}}{\xi}r^2 \\
    \therefore \ \frac{dr}{dt} \frac{1}{r} &= \frac{d\xi}{dt}\frac{1}{\xi} \\
\end{align*}

valid since $r>0$, guaranteed as long as $\xi>0$, which the $\gamma_2>0$ condition guarantees. Integrating from $0$ to $t$ gives
$$\ln r(t)-\ln r(0) = \ln\xi(t)-\ln\xi(0) \Longrightarrow r(t) = r(0)\cdot\frac{\xi(t)}{\xi(0)}$$

Applying the initial condition $r(0)=|e_d(0)|=\xi(0)$ and substituting gives
$$r(t) = \xi(0)\cdot\frac{\xi(t)}{\xi(0)} = \xi(t)\qquad\text{for all }t\ge0$$

Now that it's been show that $$\norm{e_d(t)} = \xi(t)$$ for all $t\ge0$, and $\xi(t) > 0$ since $\gamma_2>0$ then it is valid to write to write $e_d(t)$ in polar form
$$e_d(t) = \xi(t)\begin{bmatrix}\cos\theta(t) \\ \sin\theta(t)\end{bmatrix}$$
where $\theta(t)$ is the angle of the vector $e_d(t)$ in the 2D plane. Now substituting this into the ODE for $e_d(t)$ gives

\begin{align*}
    \dot{e}_d(t) &= \dot{\xi}(t)\begin{bmatrix}\cos\theta(t) \\ \sin\theta(t)\end{bmatrix} + \xi(t)\begin{bmatrix}-\sin\theta(t) \\ \cos\theta(t)\end{bmatrix}\dot{\theta}(t) \\
    &= \dot{\xi}(t)\begin{bmatrix}\cos\theta(t) \\ \sin\theta(t)\end{bmatrix} + \dot{\theta}(t) \vec{G}_1 \begin{bmatrix}\cos\theta(t) \\ \sin\theta(t)\end{bmatrix} \xi(t)
\end{align*}

Comparing coefficients from equation \eqref{eq:oscillator} shows that $\dot{\theta}(t) = \Lambda$. This indicates that $\Lambda$, which is function of $e_u, q_ew, \omega_{eu}$, determines the frequency of the oscillatory motion of $e_d(t)$.




\chapter{Zarourati Closed-loop kinematic system development}
\makeatletter\@mkboth{}{Appendix}\makeatother
\label{appen:derivations_}

$$\dot e_u = 0.5\Big(\underbrace{e_0\omega_{eu} - e_2\omega_{d3} + e_3\omega_{d2}}_{\text{this is }U\text{ from (11a)}} - k_ae_0,\vec e_a^T E\vec e_a + \vec e_a^TE(\kappa_1E\vec e_d+\kappa_2\vec e_d)\Big)$$

% This is another appendix.